\documentclass[11pt]{article}

\usepackage[margin=1in]{geometry}
\usepackage{natbib}
\usepackage{graphicx}
\usepackage{hyperref}

\title{SAE-Faithful: Are Sparse Autoencoder Features Causally Faithful to Their Auto-Generated Labels?}
\author{Aditi Shukla \and Dhruv Rai \and Simbanegavi Simbarashe}
\date{\today}

\begin{document}
\maketitle

% ---------------------------------------------------------------------
% This is a working skeleton, committed now so the document exists and
% gets added to incrementally (Instructor Review #2: "A document that
% exists gets added to; a document that does not exist gets written in a
% panic in December."). Literature review below is a first-pass draft
% covering the required >=10 cited works -- expand/refine as the semester
% continues, but the structure and citations should not need to be
% rebuilt from scratch.
% ---------------------------------------------------------------------

\begin{abstract}
Sparse Autoencoders (SAEs) decompose language-model activations into
sparse, often human-interpretable features. We investigate whether SAE
features that receive plausible-sounding auto-generated labels are also
\emph{causally faithful} to those labels -- that is, whether deliberately
strengthening a feature moves model behavior toward the concept its label
describes. We benchmark this with a steering-based evaluation across
multiple SAE architectures (ReLU, TopK, JumpReLU) and models (GPT-2 Small,
Gemma-2-2B).
\end{abstract}

\section{Introduction}
% TODO: expand from README.md Project Overview / Problem Statement.

\section{Related Work / Literature Review}
\label{sec:litreview}

This section summarizes prior work most relevant to SAE-Faithful and
states explicitly how our project differs from each.

\subsection{Sparse autoencoders for interpretability}

\citet{cunningham2023sparse} and \citet{bricken2023monosemanticity} showed
that sparse autoencoders trained on language-model activations recover
directions that are substantially more monosemantic than raw neurons,
establishing SAEs as a practical interpretability tool and introducing the
now-standard ReLU SAE formulation we use as one of our three architectures.
\emph{Difference:} both works focus on discovering and qualitatively
describing interpretable features; neither tests whether a feature's label
is \textbf{causally} accurate via steering, which is our central benchmark.

\citet{gao2024scaling} introduced the TopK SAE, which enforces a fixed
number of active features per example instead of relying on an L1 penalty,
improving the accuracy/sparsity tradeoff at scale. \citet{rajamanoharan2024jumping}
introduced JumpReLU SAEs, which use a learned per-feature activation
threshold to improve reconstruction fidelity without sacrificing sparsity.
\emph{Difference:} both papers evaluate architecture quality via
reconstruction loss and downstream-task probing, not via a causal
steering-based faithfulness benchmark comparing labels to behavior -- our
Phase 6 cross-architecture comparison is designed to test whether these
architectural differences translate into different causal faithfulness,
not just different reconstruction quality.

\subsection{Base models and pretrained SAE resources}

\citet{radford2019gpt2} introduced GPT-2, the primary model family (GPT-2
Small) used for our experiments, chosen for its small size and wide
availability of pretrained SAEs. \citet{gemmateam2024gemma2} introduced
Gemma 2, and \citet{gemmascope2024} released Gemma Scope, a suite of
pretrained SAEs across Gemma 2 layers; we use Gemma-2-2B with Gemma Scope
as our secondary, larger-scale comparison model (Research Question 9:
whether GPT-2-Small findings generalize). \emph{Difference:} these are
model/tooling releases rather than faithfulness studies; we build directly
on top of them rather than replicating their contributions.

\citet{saelens2024} (SAELens) and \citet{transformerlens2022}
(TransformerLens) are the libraries used throughout this project to load
models, cache activations, and register the forward hooks used for causal
steering (see \texttt{src/models/gpt2\_sae.py}, \texttt{src/sae/steering.py}).
\citet{neuronpedia} provides independently auto-generated feature labels
that we cross-check our own feature-label findings against (e.g. Feature
974 on GPT-2 Small being independently described as relating to
``mentions of Paris,'' see \texttt{results/phase1\_observations.md}).
\emph{Difference:} we treat Neuronpedia's labels as an independent
cross-check on labels, not as ground truth for causal faithfulness, which
is exactly the gap this project measures.

\subsection{Concept detection and representation-level evaluation}

\citet{marks2023geometry} (Geometry of Truth) show that true/false
statement representations in language models are approximately linearly
separable, and release a dataset of true/false statements we use for
Phase 5 concept-detection experiments (see
\texttt{cookbooks/03\_geometry\_of\_truth\_eda.ipynb}).
\citet{zou2023representation} (Representation Engineering / RepE) propose
extracting behavioral directions (e.g. honesty, harmlessness) from
contrastive example pairs and using them for both interpretation and
control. \emph{Difference:} both works focus on discovering or steering
along a \emph{model-level} representation direction; we instead ask
whether an already-labeled \emph{SAE feature} direction is faithful to its
label, using their datasets/methodology as concept-detection benchmarks
(Precision/Recall/F1/AUROC) rather than as the primary steering
mechanism.

\section{Methods}
% TODO: formalize Phases 1-7 from README.md / proposal.md into a Methods
% section (feature extraction, labeling, steering intervention, concept
% detection, cross-condition analysis).

\section{Experiments and Results}
% TODO: populate from results/phase1_observations.md,
% results/geometry_of_truth_observations.md, and results/tables/*.csv as
% each phase completes.

\section{Discussion}
% TODO.

\section{Conclusion}
% TODO.

\bibliographystyle{plainnat}
\bibliography{mybib}

\end{document}
